\todo[inline,color=gray!10]{meta: Dependencies: None.}

\section{Background}\label{background}

\label{sec:background}

\subsection{W3C Verifiable Credentials Data Model 2.0}\label{w3c-verifiable-credentials-data-model-2.0}

\label{sec:vcdm}

\todo[inline]{B1 --- Define VCDM core data model elements for Section 4. Verifiable credential, credential subject, issuer, claims, proof mechanism. Emphasize that VCDM defines structure but not domain semantics or format encoding. [Binding Claim \#2: grounding in VCDM 2.0]. Length: 5--7 sentences.}

\todo[inline,color=blue!20]{cite: W3C VCDM 2.0 specification}

\todo[inline]{B2 --- Establish format diversity and differing privacy capabilities. JSON-LD with Data Integrity Proofs, JWT-encoded VCs, AnonCreds (outside W3C but widely used). Each format has different privacy capabilities: selective disclosure, predicate proofs (ZKP). This diversity motivates the format-specific layer. Length: 3--4 sentences.}

\todo[inline,color=blue!20]{cite: W3C VC Data Integrity}

\todo[inline,color=blue!20]{cite: AnonCreds specification}

\todo[inline,color=blue!20]{cite: SD-JWT VC}

\todo[inline]{B3 (optional) --- Briefly establish governance context. EU Digital Identity Wallet (eIDAS 2.0, Architecture Reference Framework) imposes constraints on format choice, attribute disclosure, holder binding. Only include if space allows; context also appears in Sec 03. Length: 2--3 sentences.}

\todo[inline,color=blue!20]{cite: eIDAS 2.0 implementing regulation}

\todo[inline,color=blue!20]{cite: EU Architecture Reference Framework (ARF)}

\subsection{Multi-Level Modeling}\label{multi-level-modeling}

\label{sec:multi-level}

\todo[inline]{B4 --- Position three-layer metamodel within multi-level modeling literature. Atkinson \& K{\"u}hne multi-level metamodeling, linguistic vs. ontological typing, potency. Three layers (domain concept, credential schema, format-specific) as ontological levels connected by cross-level constraints --- independently governed concern spaces, not standard top-down refinement. No MDA terminology per DECISIONS.md. Length: 3--4 sentences.}

\todo[inline,color=blue!20]{cite: Atkinson \& K{\"u}hne --- multi-level modeling}

\todo[inline,color=blue!20]{cite: potency-based multi-level modeling}

\todo[inline]{B5 --- Establish specific metamodeling concepts used in Section 4. Ecore-style class diagrams as metamodel notation, instances as partial models with open/closed world semantics, cross-layer trace relationships as refinement mappings. Length: 2--3 sentences.}

\subsection{Partial Graph Modeling with Refinery}\label{partial-graph-modeling-with-refinery}

\label{sec:refinery}

\todo[inline]{B6 --- Introduce Refinery as the tool framework. Partial graph modeling framework, graph predicates as first-class constraint language, design space exploration via model generation from partial specifications. Third person (double-blind). Length: 4--5 sentences.}

\todo[inline,color=blue!20]{cite: Refinery --- partial graph modeling framework}

\todo[inline]{B7 --- Define Refinery language elements used in Section 4.4. Error predicates (constraints that must not hold), propagation rules (derived knowledge guiding generation), shadow predicates, scope constraints, partial interpretation (may/must/unknown). These are building blocks for Section 4.4 cross-layer constraints. Length: 4--5 sentences.}

\todo[inline,color=green!20]{formal: Define error predicate, propagation rule, shadow predicate, scope constraint}
